% ======================================================================
\chapter{Git LFS: arquivos grandes}
\label{cap:lfs}
% ======================================================================

\section{O problema}

O Git foi feito para texto, como código-fonte. Com arquivos grandes e binários
(vídeos, imagens pesadas, modelos de aprendizado de máquina, arquivos
\texttt{.psd} ou \texttt{.zip}), aparecem três problemas:

\begin{itemize}
  \item \textbf{Limite do GitHub:} o GitHub avisa a partir de 50 MiB e
        \textbf{recusa arquivos acima de 100 MiB} num push comum
        \cite{github_lfs}. Esse limite é do GitHub, não do Git.
  \item \textbf{Histórico que só cresce:} binários quase não se beneficiam da
        forma como o Git guarda diferenças entre versões. Cada versão nova de
        um arquivo grande ocupa quase o tamanho inteiro de novo, e todo
        \comando{git clone} baixa todas elas.
  \item \textbf{Operações lentas:} com muitos binários no histórico,
        \comando{clone}, \comando{fetch} e \comando{checkout} ficam lentos.
\end{itemize}

\section{A solução}

\begin{conceito}[Git LFS]
O \textbf{Git LFS} (\emph{Large File Storage}) é uma extensão do Git. No
repositório fica só um \textbf{ponteiro} pequeno, de texto. O conteúdo real vai
para um armazenamento separado, no próprio GitHub, e é baixado só quando
necessário. Os comandos do dia a dia (\comando{add}, \comando{commit},
\comando{push}) continuam os mesmos.
\end{conceito}

O ponteiro gravado no Git tem esta cara:

\begin{saida}
version https://git-lfs.github.com/spec/v1
oid sha256:e1ea497dccb8bd9d842c37f70699210c4db6d6981d7c41fecd1d028d16fade26
size 200000
\end{saida}

\section{Instalação}

\begin{itemize}
  \item \textbf{Windows:} já vem com o Git para Windows.
  \item \textbf{Linux (Debian e Ubuntu):} \comando{sudo apt install git-lfs}.
  \item \textbf{macOS:} \comando{brew install git-lfs}.
\end{itemize}

Depois, rode \textbf{uma vez por usuário}:

\begin{bashcode}
git lfs install
\end{bashcode}

\begin{saida}
Updated git hooks.
Git LFS initialized.
\end{saida}

Esse comando não instala o programa: ele liga o LFS na configuração do seu
Git.

\section{Usando}

\begin{bashcode}
# diz ao LFS quais arquivos ele deve cuidar
git lfs track "*.zip"
git lfs track "*.psd"

# a regra fica gravada no .gitattributes, que também vai para o repositório
git add .gitattributes
git add materiais.zip
git commit -m "chore: versiona os materiais com Git LFS"
git push

# confere o que está no LFS
git lfs ls-files
\end{bashcode}

\begin{aviso}[Configure antes de adicionar]
O \comando{git lfs track} só vale para o que for adicionado \emph{depois}. Um
arquivo grande que já entrou no histórico como arquivo comum continua lá.
Mover arquivos antigos para o LFS exige \comando{git lfs migrate}, que
\textbf{reescreve o histórico} e precisa ser combinado com a equipe.
\end{aviso}

\section{Limites e custos no GitHub}

\begin{itemize}
  \item Cada conta tem \textbf{10 GiB de armazenamento} e \textbf{10 GiB de
        banda} de LFS por mês, tanto no plano Free quanto no Pro
        \cite{github_lfs_uso}. Acima disso, o uso é cobrado.
  \item O tamanho máximo de um arquivo no LFS é de 2 GB nos planos Free e
        Pro.
  \item \textbf{Apagar um arquivo do LFS não libera espaço} sozinho. Para
        remover de verdade, é preciso reescrever o histórico ou apagar o
        repositório.
\end{itemize}

\begin{dica}[Quando não usar]
Não use o LFS para arquivos pequenos em grande quantidade nem para o que pode
ser gerado de novo (arquivos compilados, dependências): esses devem ir para o
\arq{.gitignore}.
\end{dica}
